\subsection{李代数的函子性质}

设 \(G\) 与 \(H\) 为李群，\(\phi\) 为 \(G\) 到 \(H\) 的态射。\(U(\phi)\) 在 \(L(G)\) 上的限制（它正是 \(T_e(\phi)\)）是 \(L(G)\) 到 \(L(H)\) 的连续态射，我们把它记作 \(L(\phi)\)。若 \(\psi\) 是 \(H\) 到某个李群的态射，则 \(L(\psi \circ \phi) = L(\psi) \circ L(\phi)\)。

为使 \(\phi\) 是浸入，充要条件是 \(\mathrm{L}(\phi)\) 是 \(\mathrm{L}(\mathrm{G})\) 到 \(\mathrm{L}(\mathrm{H})\) 的某个具有拓扑余的子代数上的同构。特别地，若 \(\mathrm{G}\) 是 \(\mathrm{H}\) 的李子群且 \(\phi\) 是典范单射，则 \(\mathrm{L}(\mathrm{G})\) 与 \(\mathrm{L}(\mathrm{H})\) 的一个李子代数借助 \(\mathrm{L}(\phi)\) 认同。更特殊地，若 \(\mathrm{G}\) 是 \(\mathrm{H}\) 的开子群，则 \(\mathrm{L}(\mathrm{G}) = \mathrm{L}(\mathrm{H})\)。

若 G 是 H 的李拟子群，则 L(G) 也与 L(H) 的一个闭李子群认同。

为使 \(\phi\) 是浸没，充要条件是 \(\mathrm{L}(\phi)\) 是满射的且其核具有拓扑余。这时，\(N\)（\(\phi\) 的核）是 \(G\) 的李子群且 \(\mathrm{L}(N) = \mathrm{Ker} \, L(\phi)\)。特别地，若 \(H\) 是 \(G\) 关于一个正规李子群 \(P\) 的商李群，则 \(L(P)\) 是 \(\mathrm{L}(G)\) 的理想，且若 \(\phi\) 是 \(G\) 到 \(H\) 上的典范满射，则 \(\mathrm{L}(G/P)\) 与 \(\mathrm{L}(G)/\mathrm{L}(P)\) 借助由 \(\mathrm{L}(\phi)\) 过渡到商而导出的态射认同。

设 \(I\) 为有限集，\((\mathbf{G}_{i})_{i\in I}\) 为李群族，\(G\) 为它们的积，\(p_i\) 为 \(G\) 到 \(G_i\) 上的典范态射。于是 \((\mathrm{L}(p_i))_{i\in I}\) 是李代数 \(L(G)\) 到李代数 \(\prod_{i\in I} L(G_i)\) 中的态射，且是可赋范空间之间的同构。因而 \(L(G)\) 与 \(\prod_{i\in I} L(G_i)\) 借助 \((L(p_i))_{i\in I}\) 认同。

命题 28. 设 G 与 H 为李群，\(\phi\) 为 G 到 H 的态射。设 K 的特征为 0 且 H 是有限维的。

\begin{enumerate}
\def\labelenumi{(\roman{enumi})}
\item
  \(\phi\) 的核 N 是 \(\mathbf{G}\) 的李子群，且 \(\mathrm{L}(\mathbf{N}) = \mathrm{Ker} \, \mathrm{L}(\phi)\)。
\item
  态射 \(\psi\)（由 \(\phi\) 过渡到商而导出的、G/N 到 H 中的）是浸入。
\item
  若 \(\phi(\mathbf{G})\) 在 \(\mathbf{H}\) 中是闭的且 \(\mathbf{G}\) 的拓扑具有可数基，则 \(\phi(\mathbf{G})\) 是 \(\mathbf{H}\) 的李子群，\(\psi\) 是李群 \(\mathbf{G}/\mathbf{N}\) 到李群 \(\phi(\mathbf{G})\) 上的同构，且 \(\mathbf{L}(\phi(\mathbf{G})) = \mathbf{Im} \mathbf{L}(\phi)\)。
\end{enumerate}

设 G 通过映射 \((g, h) \mapsto \phi(g)h\) 左作用在 H 上。只需对 e 的轨道引用 §1 第 7 节的命题 14。

命题 29. 设 G 与 H 为李群，\(\phi\) 为 G 到 H 的态射。设 K 的特征为 0 且 H 是有限维的。若 H' 是 H 的李子群，则 G' = \(\phi^{-1}(\mathrm{H}')\) 是 G 的李子群且 L(G') = L(φ)\(^{-1}\)(L(H'))。

设 \(\pi\) 为 H 到齐性空间 \(\mathbf{X} = \mathbf{H} / \mathbf{H}'\) 的典范映射。设 G 左作用在 \(\mathbf{X}\) 上，作用映射为 \((g, x) \mapsto \phi(g)x\)。\(\pi(e)\) 的稳定子群是 \(G'\)，因而它是 G 的李子群 (§ 1, 第 7 节, 命题 14)。\(\pi(e)\) 的轨道映射是 \(\pi \circ \phi\)。由 § 1 的命题 14，L(G') 是 L(\(\pi \circ \phi\)) = \(T_e(\pi)\) \(\circ\) L(\(\phi\)) 的核。\(T_e(\pi)\) 的核是 L(H') (§ 1, 第 6 节, 命题 11 (i))，因而 Ker L(\(\pi \circ \phi\)) = L(\(\phi\)) \(^{-1}\)(L(H'))。

推论 1. 设 \(\mathbf{G}\)、\(\mathbf{H}\) 为李群，\(\phi_1\) 与 \(\phi_2\) 为 \(\mathbf{G}\) 到 \(\mathbf{H}\) 的态射。设 \(\mathbf{K}\) 的特征为 0 且 \(\mathbf{H}\) 是有限维的。由 \(g \in \mathbf{G}\)（满足 \(\phi_1(g) = \phi_2(g)\)）组成的集合是李子群 \(\mathbf{G}'\)（\(\mathbf{G}\) 的李子群），且 \(\mathbf{L}(\mathbf{G}')\) 是由 \(x \in \mathbf{L}(\mathbf{G})\)（满足 \(\mathrm{L}(\phi_1)x = \mathrm{L}(\phi_2)x\)）组成的集合。

我们记 \(\phi(g) = (\phi_1(g), \phi_2(g))\)（对所有 \(g \in G\)），于是 \(\phi\) 是 \(G\) 到 \(H \times H\) 的态射。设 \(\Delta\) 为 \(H \times H\) 的对角子群。于是 \(G' = \phi^{-1}(\Delta)\)，且 \(L(\phi)x = (L(\phi_1)x, L(\phi_2)x)\)（对所有 \(x \in L(G)\)）。现在只需引用命题 29。

推论 2. 设 \(\mathbf{G}\) 为有限维李群，\(\mathbf{G}_1\) 与 \(\mathbf{G}_2\) 为 \(\mathbf{G}\) 的两个李子群。设 \(\mathbf{K}\) 的特征为 0。则 \(G_1 \cap G_2\) 是 \(\mathbf{G}\) 的李子群，其李代数为 \(\mathbf{L}(\mathbf{G}_1) \cap \mathbf{L}(\mathbf{G}_2)\)。

我们对 \(\mathbf{G}_1\) 到 \(\mathbf{G}\) 的典范单射与子群 \(\mathbf{G}_2\) 应用命题 29。

推论 3. 设 G、G'、H 为李群，\(\phi: G \to H\) 与 \(\phi': G' \to H\) 为李群态射。设 K 的特征为 0 且 H 是有限维的。设 F 为由 \((g, g') \in G \times G'\)（满足 \(\phi(g) = \phi'(g')\)）组成的集合。则 F 是 \(G \times G'\) 的李子群，且 L(F) 是由 \((x, x') \in L(G) \times L(G')\)（满足 L( \(\phi\) ) \(x = L(\phi')x'\)）组成的集合。

我们对态射 \((g, g') \mapsto \phi(g)\) 与 \((g, g') \mapsto \phi'(g')\)（从 \(G \times G'\) 到 \(H\) 中）应用推论 1。

命题 30. 设 G 为具有可数基的有限维李群，且

H 与 H' 为 G 的李子群。设 K 的特征为 0 且 HH' 在 G 中是局部闭的。

\begin{enumerate}
\def\labelenumi{(\roman{enumi})}
\item
  HH' 是 G 的子流形，且 \(\mathrm{T}_{e}(\mathrm{HH}') = \mathrm{L}(\mathrm{H}) + \mathrm{L}(\mathrm{H}')\)。
\item
  设 H 的每个元素与 H' 的每个元素交换。则 HH' 是 G 的李子群。设 \(\phi\) 为映射 \((h, h') \mapsto hh'\)（从 H \(\times\) H' 到 HH' 上）。\(\phi\) 的核是对 \((m, m^{-1})\)（其中 \(m \in \mathrm{H} \cap \mathrm{H}'\)）组成的集合，且 \((\mathrm{H} \times \mathrm{H}') / \mathrm{Ker} \phi\) 到 HH' 上的、由 \(\phi\) 过渡到商而导出的态射是李群同构。
\end{enumerate}

设 \(\mathbf{H} \times \mathbf{H}'\) 右作用在 \(\mathbf{G}\) 上，作用映射为 \(((h, h'), g) \mapsto hgh'^{-1}\)。设 \(\rho\) 为 \(e\) 的轨道映射：\((h, h') \mapsto hh'^{-1}\)。由 § 1 第 7 节的命题 14 (iii)，HH' 是 \(\mathbf{G}\) 的子流形且 \(\mathrm{T}_e(\mathrm{HH}') = \mathrm{Im} \mathrm{T}_e(\rho)\)。现在

\[
\mathrm{T} _ {e} (\rho) (\mathrm{L} (\mathrm{H}) \times \{0 \}) = \mathrm{L} (\mathrm{H}) \quad \text { and } \quad \mathrm{T} _ {e} (\rho) (\{0 \} \times \mathrm{L} (\mathrm{H} ^ {\prime})) = \mathrm{L} (\mathrm{H} ^ {\prime})
\]

因而 \(\mathrm{T}_{e}(\mathrm{HH}^{\prime})=\mathrm{L}(\mathrm{H})+\mathrm{L}(\mathrm{H}^{\prime})\)。设 H 的每个元素与 \(H^{\prime}\) 的每个元素交换。则 \(HH^{\prime}\) 是 G 的子群。由 (i)，它是 G 的李子群。断言的其余部分由命题 28 得出。

命题 31. 设 G 为具有可数基的有限维李群，H 为 G 的正规李子群，A 为 G 的李子群。设 K 的特征为 0 且 AH 是闭的。设 \(\phi\) 为 G 到 G/H 上的典范态射。则典范映射

\[
\mathrm{A} / (\mathrm{H} \cap \mathrm{A}) \rightarrow \phi (\mathrm{A}), \quad \mathrm{AH} / \mathrm{H} \rightarrow \phi (\mathrm{A})
\]

都是李群同构。

由命题 30，AH 是 G 的李子群。由命题 29 的推论 2，H ∩ A 是 G 的李子群。因此谈论群 AH/H 与 A/(H ∩ A) 是有意义的。另一方面，φ(A) 是 AH 在 G/H 中的典范像，它是闭的，因而是 G/H 的李子群（命题 28 (iii)）。把命题 28 应用于复合态射 A → G → G/H 与 AH → G → G/H，便证明了命题中的典范映射是李群同构。

命题 32. 设 G 与 H 为李群，k 为 K 的非离散闭子域，\(\phi\) 为 \(k\) 上的李群结构之间 G 到 H 的态射。设 K 的特征为 0。若 \(L(\phi)\) 是 K-线性的，则 \(\phi\) 是 K 上的李群结构之间的态射。

对所有 \(g \in G\)，

\[
\mathrm{T} _ {g} (\phi) = \mathrm{T} _ {e} (\gamma (\phi (g))) \circ \mathrm{L} (\phi) \circ \mathrm{T} _ {g} (\gamma (g) ^ {- 1})
\]

因而 \(\mathrm{T}_g(\phi)\) 是 K-线性的。命题于是由 Differentiable and Analytic Manifolds, R, 5.14.6 得出。

